% SSSUP-FONTE-STORICA: fonti-locali/originali/2014-2020.zip
% Anno 2017, esercizio 3. Testo storico; figura TikZ ricostruita in PSTricks.
% Verificare editorialmente la trasposizione prima di considerarla definitiva.
\renewcommand{\Anno}{2017}
\renewcommand{\NumeroProblema}{3}

\renewcommand{\Testo}{
Una scacchiera $n\times m$ (in figura quella $3\times5$) contiene $nm-1$ piastrelle quadrate ed uno spazio vuoto (in nero nella figura). Una {\it mossa} è lo spostamento di una piastrella in un quadrato adiacente (che deve essere necessariamente vuoto).
\begin{center}
% SSSUP-FIGURA-RICOSTRUITA: scacchiera 3 x 5, originale TikZ 2017.
\psset{unit=0.57cm}
\begin{pspicture}(-0.65,-0.48)(5.40,3.70)
  \multido{\n=0+1}{6}{\psline[linecolor=gray,linewidth=.45pt](\n,0)(\n,3)}
  \multido{\n=0+1}{4}{\psline[linecolor=gray,linewidth=.45pt](0,\n)(5,\n)}
  \psframe[fillstyle=solid,fillcolor=black,linestyle=none](2,1)(3,2)
  \rput(0.5,3.40){1}
  \rput(1.5,3.40){2}
  \rput(2.5,3.40){3}
  \rput(3.5,3.40){4}
  \rput(4.5,3.40){5}
  \rput(-.36,2.5){1}
  \rput(-.36,1.5){2}
  \rput(-.36,0.5){3}
\end{pspicture}
\end{center}

\begin{enumerate}
\item[(a)] Mostrare che per spostare lo spazio vuoto dalla posizione $(h,k)$ alla posizione $(h+p, k+q)$ il numero minimo di {\it mosse} è $|p|+|q|$.

\item[(b)] Nel caso di scacchiera $3\times5$ determinare il minimo numero di {\it mosse} per portare una piastrella dalla posizione $(1,1)$ alla posizione $(3,5)$ nel caso in cui lo spazio vuoto si trovi inizialmente in $(2,1)$ e nel caso in cui lo spazio vuoto si trovi inizialmente in $(1,2)$.

\item[(c)] Nel caso di scacchiera $3\times5$, sapendo che lo spazio vuoto potrebbe inizialmente trovarsi in una qualsiasi casella diversa da $(1,1)$, determinare il {\it Valor Medio} del numero minimo di {\it mosse} per spostare una piastrella dalla posizione $(1,1)$ alla posizione $(3,5)$.
\end{enumerate}
}

\renewcommand{\Risposta}{
% SSSUP-BOZZA-AI: Risposta; autore=ChatGPT; data=2026-11-10; stato=DA-VERIFICARE
% SSSUP-PROVENIENZA: ChatGPT (OpenAI), sessione ProblemiSSSUP del 2026-11-10; derivato dalla traccia canonica; revisione umana pendente.
(a) Il minimo \`e $|p|+|q|$, purch\'e le due caselle indicate
appartengano alla scacchiera.

(b) Occorrono $20$ mosse se il vuoto parte da $(2,1)$,
e $18$ mosse se parte da $(1,2)$.

(c) Se il vuoto parte da $(r,c)\ne(1,1)$, il minimo \`e
\[
\boxed{M(r,c)=18+|r-1|+|c-2|.}
\]
La media aritmetica sulle $14$ posizioni possibili \`e
\[
\boxed{\frac1{14}\sum M(r,c)=\frac{41}{2}=20.5.}
\]
}
\renewcommand{\Aiutino}{
% SSSUP-BOZZA-AI: Aiutino; autore=ChatGPT; data=2026-11-10; stato=DA-VERIFICARE
% SSSUP-PROVENIENZA: ChatGPT (OpenAI), sessione ProblemiSSSUP del 2026-11-10; derivato dalla traccia canonica; revisione umana pendente.
Per (a), ogni mossa del vuoto cambia una sola coordinata
di una unit\`a.
Per (b) e (c), seguire soltanto la \emph{piastrella designata}
e la posizione del vuoto: delle altre piastrelle non serve conoscere
l'identit\`a. Quando la piastrella designata avanza di una casella,
lo spazio vuoto resta nella sua posizione precedente.
Quante mosse servono per portare di nuovo il vuoto davanti
alla piastrella, a seconda che questa continui diritta
o cambi direzione?
}
\renewcommand{\SoluzioneUfficiale}{}
\renewcommand{\SoluzioneRedazionale}{
% SSSUP-BOZZA-AI: SoluzioneRedazionale; autore=ChatGPT; data=2026-11-10; stato=DA-VERIFICARE
% SSSUP-PROVENIENZA: ChatGPT (OpenAI), sessione ProblemiSSSUP del 2026-11-10; derivato dalla traccia canonica; revisione umana pendente.
\textbf{(a) Spostamento dello spazio vuoto.}
Ogni mossa scambia lo spazio vuoto con una piastrella adiacente,
e quindi ne modifica una sola coordinata di una unit\`a.
Per passare da $(h,k)$ a $(h+p,k+q)$ servono almeno $|p|+|q|$ mosse.
Viceversa, spostando il vuoto lungo un cammino formato da
$|p|$ passi verticali e $|q|$ orizzontali, senza uscire dal rettangolo,
si raggiunge la destinazione in esattamente tante mosse.

\textbf{(b) Piastrella designata.}
Seguiamo soltanto la piastrella che inizialmente si trova
in $(1,1)$, che chiameremo \emph{piastrella speciale},
e il vuoto. Ogni spostamento della piastrella speciale \`e una mossa,
mentre il vuoto pu\`o anche muoversi scambiandosi con le altre piastrelle.

Indichiamo con $k$ il numero di spostamenti della piastrella speciale.
Poich\'e deve andare da $(1,1)$ a $(3,5)$,
la distanza di Manhattan fra origine e destinazione \`e $6$.
Dunque $k\ge6$ e $k$ \`e pari. Possiamo inoltre supporre che
due spostamenti consecutivi della piastrella speciale
non si annullino a vicenda: se la piastrella passasse da $A$ a $B$
e subito dopo da $B$ ad $A$, la configurazione formata da
\emph{piastrella speciale e vuoto} tornerebbe identica a prima
dei due scambi, quindi potremmo eliminare l'intero tratto intermedio.

Fra due spostamenti consecutivi che non si annullano,
il vuoto deve passare dalla casella appena lasciata dalla
piastrella alla casella in cui questa si sposter\`a.
Se la piastrella \emph{cambia direzione}, tali due caselle
sono a distanza di Manhattan $2$, e servono almeno due
mosse del vuoto. Se la piastrella procede \emph{diritta},
le due caselle stanno ai lati opposti della posizione attuale
della piastrella: il vuoto deve aggirarla, e servono
almeno quattro mosse.

Se $k\ge8$, anche usando soltanto due mosse del vuoto
fra tutti gli spostamenti consecutivi della piastrella speciale,
si ottiene il limite inferiore
\[
 k+2(k-1)\ge22. \tag{1}
\]
Per trovare il minimo basta quindi considerare $k=6$,
se \`e disponibile una costruzione con non pi\`u di $20$ mosse.
In tal caso la piastrella speciale deve fare esattamente
quattro passi a destra (D) e due in basso (B), in qualche ordine.

Se il primo passo \`e a destra, la sequenza di sei passi
pu\`o cambiare direzione al massimo quattro volte.
Nei cinque passaggi fra uno spostamento della piastrella
e il successivo ve ne sar\`a dunque almeno uno diritto:
servono almeno
\[
6+4\cdot2+1\cdot4=18\quad\text{mosse}. \tag{2}
\]
Se invece il primo passo \`e in basso, la sequenza pu\`o
cambiare direzione al massimo tre volte; ne segue
\[
6+3\cdot2+2\cdot4=20\quad\text{mosse}. \tag{3}
\]
Questi limiti non contano eventuali mosse \emph{iniziali}
necessarie per portare il vuoto a contatto con la piastrella speciale.

Il limite $18$ \`e raggiungibile con il vuoto inizialmente
in $(1,2)$. Basta far seguire alla piastrella speciale
la sequenza di passi
\[
 \text{destra, basso, destra, basso, destra, destra}.
\]
Fra gli spostamenti della piastrella, il vuoto si porta
davanti a essa con due mosse per ciascuno dei quattro
cambi di direzione e con quattro mosse nell'unico tratto diritto.
Tutti questi piccoli aggiramenti si possono effettuare
all'interno della scacchiera $3\times5$, dando in totale $18$ mosse.

Se il vuoto parte in $(2,1)$, un primo passo della
piastrella speciale verso il basso richiede almeno $20$ mosse,
per (3); per un primo passo verso destra, occorrono
almeno due mosse preliminari per condurre il vuoto da
$(2,1)$ a $(1,2)$, e poi almeno $18$ mosse, per (2).
Il risultato \`e quindi ancora $20$, ed \`e raggiungibile
portando prima il vuoto in $(1,2)$ attraverso $(2,2)$.
Il limite (1) esclude percorsi pi\`u lunghi che possano migliorare
questi valori.

\textbf{(c) Media sulle posizioni iniziali del vuoto.}
Sia $(r,c)$ una delle $14$ caselle diverse da $(1,1)$.
Prima del primo spostamento della piastrella speciale,
il vuoto deve raggiungere una delle due caselle adiacenti
$(1,2)$ e $(2,1)$. Poniamo
\[
 d_R=|r-1|+|c-2|,\qquad d_B=|r-2|+|c-1|.
\]
I cammini pi\`u brevi del vuoto verso $(1,2)$ evitano
sempre la casella $(1,1)$, dunque richiedono esattamente
$d_R$ mosse. Dalle stime precedenti, se la piastrella speciale
compie sei passi, il numero totale delle mosse \`e almeno
\[
 \min\{d_R+18,\ d_B+20\}=d_R+18,
\]
poich\'e la disuguaglianza triangolare d\`a $d_R\le d_B+2$.
Se invece la piastrella speciale compie almeno otto passi,
dalle stesse osservazioni e da (1) servono almeno
\[
\min\{d_R,d_B\}+22\ge d_R+20>d_R+18
\]
mosse. Portando prima il vuoto in $(1,2)$ lungo un cammino
minimo e poi eseguendo la costruzione di $18$ mosse,
si ottiene esattamente $d_R+18$. Quindi
\[
 \boxed{M(r,c)=18+|r-1|+|c-2|.}
\]

I valori minimi, con righe e colonne corrispondenti alle caselle
della scacchiera, sono
\[
\begin{array}{c|ccccc}
 &1&2&3&4&5\\ \hline
1&\text{--}&18&19&20&21\\
2&20&19&20&21&22\\
3&21&20&21&22&23
\end{array}
\]
La somma delle distanze $d_R$ sulle $14$ caselle \`e
$6+12+17=35$. Interpretando il ``valor medio'' richiesto
come media aritmetica sulle $14$ posizioni iniziali,
si ottiene
\[
\frac1{14}\sum_{(r,c)\ne(1,1)} M(r,c)
=18+\frac{35}{14}=\boxed{\frac{41}{2}}.
\]

\textit{Nota sul testo.} Se si interpreta ``valor medio'' come
valore atteso, \`e necessario supporre che tutte le posizioni
iniziali del vuoto siano equiprobabili; il testo originale
non esplicita questa ipotesi.
}
\renewcommand{\SoluzioniAlternative}{}
\renewcommand{\Osservazioni}{}
\renewcommand{\Commenti}{}
